% Solution: Cabral (2017) Ch. 9 Exercise 9.15 "Multimarket contact". Input by \solitem. Check: check_T3.py (item 21).
\begin{exsol}{Cabral Exercise 9.15}{244}{Advanced}{4}
\asked Two markets with the same inelastic demand ($q$ units up to a reservation price $p^M$). Firm~1's
cost is $\underline c$ in Market~1 and $\overline c$ in Market~2, and the reverse for Firm~2, with
$\underline c < \overline c < p^M$. Firms set prices in both markets each period. Find the minimum discount factor
for which the optimal collusive solution is stable.

\method Cabral \S9.3 (pp.~227--231, multimarket contact). The efficient collusive agreement in each market is
that the low-cost firm sets $p^M$ and the high-cost firm sets a higher price and sells nothing. The punishment is a
price war at the static equilibrium, in which the low-cost firm prices just below the rival's cost
$\overline c$.

\solution
\textbf{Collusion.} Each firm sells all $q$ units at $p^M$ in its low-cost market: profit per period
$A = q(p^M - \underline c)$.

\textbf{Deviation.} Take Firm~1. In Market~1 it is already a monopolist, and in Market~2 it can undercut Firm~2's price
$p^M$ slightly, gaining the whole market at cost $\overline c$: extra profit $B = q(p^M - \overline c) > 0$.
So the deviation profit in the period is $A + B$.

\textbf{Punishment.} Static equilibrium in each market: the low-cost firm charges just below $\overline c$ and
the high-cost firm sells nothing. Each firm earns, in its low-cost market, $q(\overline c - \underline c) = A - B$ per period.
Both markets are in the war.

\textbf{Condition.}
\[
  \frac{A}{1-\delta} \ge A + B + \frac{\delta (A - B)}{1-\delta}
  \iff A \ge (A + B)(1-\delta) + \delta(A - B)
  \iff 0 \ge B - 2\delta B
  \iff \delta \ge \frac12 .
\]
So the minimum discount factor is $\delta = \frac12$, whatever $\underline c$, $\overline c$, $p^M$ or $q$ (with
$\overline c < p^M$ so that $B > 0$).

\interp In a single market the efficient agreement is never stable: the high-cost firm earns zero and would
undercut at any $\delta$, since no punishment can hurt it. With two markets each firm is the ``low-cost'' firm in one
market and has something to lose in the other, so the threat of retaliation in the rival's market holds the agreement
together. This is Cabral's point that multimarket contact matters when markets are asymmetric
(his airlines and dog food boxes, pp.~229--230).

\verify \texttt{check\_T3.py} (item~21) solves the indifference equation symbolically with $A$, $B$ and
$A - B$ in terms of $p^M$, $\underline c$, $\overline c$, $q$ and obtains $\frac12$. It tests numbers
($p^M = 10$, $\underline c = 2$, $\overline c = 6$, $q = 1$: $A = 8$, $B = 4$, war profit $4$) at $\delta = \frac12,
\frac23$ (stable) and $0.49$ (not stable). Cabral prints no answer.

\mistakes
\begin{itemize}
  \item Using zero as the punishment payoff. In the static equilibrium of this asymmetric game the low-cost firm
  still earns $q(\overline c - \underline c)$ in its own market.
  \item Counting only one market's deviation gain: Firm~1 deviates in Market~2 (the rival's market) and
  continues to collect $A$ in Market~1 in the same period.
\end{itemize}

\answer{$\delta \ge \dfrac12$: with collusive profit $A$, extra deviation profit $B$ and war profit $A - B$ the
condition $\frac{A}{1-\delta} \ge A + B + \frac{\delta(A-B)}{1-\delta}$ reduces to $\delta \ge \frac12$, independent
of the costs, $q$ and $p^M$.}
\end{exsol}
